\chapter{Introduccion}

\section{Pregunta motora}

La pregunta de fondo de esta investigacion es si la estructura covariante
relativista en $3+1$ dimensiones --- algebra de Clifford, generadores de
Lorentz, $SL(2,\mathbb{C})$, fermiones de Dirac y Weyl --- puede entenderse
como un \emph{sector efectivo infrarrojo} de un sustrato discreto local mas
fundamental, al que llamamos \emph{protoespacio}.

\section{Hipotesis de trabajo}

\begin{enumerate}[label=(H\arabic*)]
  \item \textbf{Emergencia infrarroja:} el continuo relativista no describe el
        objeto microscopico exacto, sino solo su sector de baja energia y
        longitudes de onda grandes.
  \item \textbf{Primacia de la localidad:} la propiedad robusta de partida es
        la localidad del paso unitario \Uop. La covariancia Lorentz aparece
        despues, si aparece.
  \item \textbf{Sustrato ampliado:} el protoespacio no es la red espacial
        desnuda, sino la tupla
        \[\Proto \sim \ProtoTuple.\]
\end{enumerate}

\section{Plan del libro}

La Parte~I construye la escalera $1D \to 2D \to 3D \to 3+1$ a traves de SSH,
grafeno, semimetales de Weyl/Dirac y caminatas cuanticas. La Parte~II extrae
la estructura espinorial y los generadores de Lorentz como objetos
emergentes. La Parte~III enuncia criterios para que un sistema discreto
admita un cono de luz, isotropia e invariancia Lorentz efectivos. La
Parte~IV deja abiertos los dos nudos identificados en la revision del 28 de
marzo: el espectro global tipo Brillouin del paso discreto, y la
formulacion del protoespacio independiente del lenguaje reciproco.
